% !TeX root = ../Tex/main.tex

\begin{center}
    \subsection*{Abstract}
\end{center}
This article explores the topic of linguistic simplification in parliaments. In 1868 a law broadened the electorate base of UK House of Commons. Can we exploit this setting to measure the effects on linguistic simpilification?
\textcite{spirlingDemocratizationLinguisticComplexity2016} proved that a broader electoral base has different effects on linguistic simpilification for cabinet and non cabinet members.
Building on  \textcite{spirlingDemocratizationLinguisticComplexity2016}, I try to support the causal inference analysis with further investigation of data. In particular, I disentagle the effects of to different measures on the outcome variable provided by \textcite{spirlingDemocratizationLinguisticComplexity2016}. I will also reduce the time window to empower the effects of observations around 1868 cutoff.

\vspace*{1em}
Text as data techinques are used the selct the best method to get a DTM matrix out of the raw speeches. Causal Inference methods are employed on this data. After these steps, I lay down the data generating process, I use a causal diaciclycal graph to build the causal mechanism. These sections should be a suitable theoretical ground for the application of an regression discontinuity design.

%\lipsum[1] % genera il primo paragrafo di testo Lorem Ipsum

%------- End of the File -------